\documentclass[a4paper,11pt]{article}

\usepackage[utf8]{inputenc}
\usepackage[T1]{fontenc}
\usepackage{listings}
\usepackage{xcolor}
\usepackage{hyperref}
\usepackage{geometry}
\usepackage{amssymb,amsmath}
\usepackage{enumitem}
\usepackage{tabularx}
\usepackage{fancyvrb}

\geometry{margin=2.5cm}

\definecolor{javared}{RGB}{180,0,0}
\definecolor{javagreen}{RGB}{0,128,0}
\definecolor{javapurple}{RGB}{128,0,128}
\definecolor{javacomment}{RGB}{128,128,128}

\lstdefinelanguage{JavaFiltered}{
  morekeywords={abstract,assert,boolean,break,byte,case,catch,char,class,
    const,continue,default,do,double,else,enum,extends,false,final,finally,
    float,for,goto,if,implements,import,instanceof,int,interface,long,native,
    new,null,package,private,protected,public,return,short,static,strictfp,
    super,switch,synchronized,this,throw,throws,transient,true,try,void,
    volatile,while,var},
  morecomment=[l]{//},
  morecomment=[s]{/*}{*/},
  morestring=[b]",
  sensitive=true
}

\lstset{
  language=JavaFiltered,
  basicstyle=\small\ttfamily,
  keywordstyle=\color{javapurple}\bfseries,
  commentstyle=\color{javacomment},
  stringstyle=\color{javared},
  numbers=left,
  numberstyle=\tiny\color{javacomment},
  frame=single,
  breaklines=true,
  tabsize=2,
  columns=fullflexible
}

\title{\textbf{PadePrinter: A Complete Line-by-Line Walkthrough}\\[4pt]
       \large arb4j --- M\"untz--Pad\'e Spectral-Tau Printer}
\author{Stephen Crowley}
\date{\today}

\begin{document}
\maketitle
\tableofcontents
\newpage

\section{Introduction}

PadePrinter is a \texttt{picocli}-based command-line application that prints the
Pad\'e approximant and Christoffel--Darboux kernel of the M\"untz--Pad\'e
spectral-tau solution of the constant-coefficient fractional Riccati equation:

\[
D^{\mu} y(v,t) = P(v) + Q(v)\cdot y(v,t) + R(v)\cdot y(v,t)^{2},\qquad y(0)=0.
\]

It lives at \texttt{src/main/java/arb/applications/PadePrinter.java} and
implements \texttt{Runnable} so that \texttt{picocli} can dispatch to it.

\section{Package and Imports (Lines 1--20)}

\begin{lstlisting}
package arb.applications;

import java.io.InputStream;
import java.util.Properties;

import arb.*;
import arb.Integer;
import arb.expressions.Context;
import arb.functions.complex.ComplexFunction;
import arb.functions.complex.MuntzPadeApproximant;
import arb.functions.complex.RiccatiMuntzPadeFunctional;
import arb.functions.integer.ComplexPolynomialSequence;
import arb.functions.integer.ComplexSequence;
import arb.functions.real.RealFunction;
import arb.stochastic.processes.heston.RoughHestonRiccatiMuntzPadeFunctional;
import arb.utensils.ShellFunctions;
import arb.utensils.text.tables.TextTable;
import picocli.CommandLine;
import picocli.CommandLine.Command;
import picocli.CommandLine.Option;
\end{lstlisting}

\textbf{Line 1:} Package declaration. PadePrinter is an application, not a
library component, so it lives in \texttt{arb.applications}.

\textbf{Lines 3--4:} \texttt{InputStream} and \texttt{Properties}---used by the
\texttt{VersionProvider} inner class to read build metadata from
\texttt{/arb/build-info.properties} baked into the jar.

\textbf{Line 6:} Wildcard import of the \texttt{arb} package, pulling in
\texttt{Real}, \texttt{Complex}, \texttt{ComplexPolynomial}, and every other
core type.

\textbf{Line 7:} Explicit import of \texttt{arb.Integer} (the ball-arithmetic
variant, not \texttt{java.lang.Integer})---necessary because the wildcard
import can shadow it.

\textbf{Lines 8--17:} The domain types. \texttt{Context} wires shared variables
and sub-functions. \texttt{ComplexFunction} is the expression-compiler's
function interface for complex codomains.
\texttt{MuntzPadeApproximant} is the result of the Riccati solver---it holds
every computed sequence ($\sigma$-table, Jacobi coefficients, orthogonal
polynomials, kernel). \texttt{RiccatiMuntzPadeFunctional} defines the Riccati
equation with coefficients $P$, $Q$, $R$. \texttt{RoughHestonRiccatiMuntzPadeFunctional}
specialises it for the rough Heston model. \texttt{ComplexPolynomialSequence}
and \texttt{ComplexSequence} are the expression-compiler's typed sequence
interfaces. \texttt{RealFunction} wraps a real-valued compiled expression.
\texttt{ShellFunctions} provides the \texttt{plot()} entry point for JavaFX
plotting. \texttt{TextTable} formats tabular data to stdout.

\textbf{Lines 18--20:} \texttt{picocli} annotations and entry point. The
\texttt{@Command}, \texttt{@Option} annotations define the CLI interface;
\texttt{CommandLine} executes it.

\section{Class Javadoc and \texttt{@Command} Annotation (Lines 22--107)}

\begin{lstlisting}[firstnumber=22]
/**
 * Prints the first N orthogonal polynomials of the numerator and denominator of
 * the M\"untz-Pad\'e spectral-tau solution of the constant-coefficient
 * (time-independent) fractional Riccati equation.
 *
 * <pre>
 *   D^\mu y(v,t) = P(v) + Q(v)\cdot y(v,t) + R(v)\cdot y(t,v)^2
 *   y(0) = 0
 * </pre>
 *
 * @author Stephen Crowley \copyright2024--2026
 * @see arb.documentation.BusinessSourceLicenseVersionOnePointOne \copyright terms
 */
\end{lstlisting}

\textbf{Lines 22--34:} The class-level Javadoc states the problem. The equation
here shows $y$ depending on both $v$ (a perturbation parameter) and $t$. The
M\"untz--Pad\'e method solves this by expanding $y$ as a M\"untz series in
$z = t^{\mu}$.

\newpage

\begin{lstlisting}[firstnumber=35]
@Command(name = "pade", description =
{ "",
  "Print the Pad\'e approximant and Christoffel--Darboux kernel",
  ...
  "" }, versionProvider = PadePrinter.VersionProvider.class,
        mixinStandardHelpOptions = true)
\end{lstlisting}

\textbf{Line 35:} The \texttt{@Command} annotation registers the CLI as
\texttt{pade}. \texttt{mixinStandardHelpOptions = true} auto-adds
\texttt{--help} and \texttt{--version} flags.

\textbf{Lines 36--107:} The \texttt{description} array is a multi-line string
that \texttt{picocli} prints when \texttt{--help} is invoked. It documents:
\begin{itemize}[nosep]
\item The Riccati recurrence for the coefficients $a_k$
\item The moment functional $\mu(z^k) = m(k) = a(k+1)$
\item The $\sigma$-table as a Lanczos/Danilevsky reduction of the Hankel
  moment matrix
\item The termination theorem: when $0 \in \Delta\beta_n$, the Jacobi
  recurrence has converged
\item The Christoffel--Darboux kernel and its formula
\item The discriminant $\Delta = Q^2 - 4PR$ and its role
\end{itemize}

\section{\texttt{VersionProvider} Inner Class (Lines 111--135)}

\begin{lstlisting}[firstnumber=111]
static class VersionProvider implements
                                 picocli.CommandLine.IVersionProvider
  {
    @Override
    public String[] getVersion() throws Exception
    {
      Properties props = new Properties();
      try ( InputStream in = PadePrinter.class
              .getResourceAsStream("/arb/build-info.properties"))
      {
        if (in != null)
        {
          props.load(in);
        }
      }
      String arb4j = props.getProperty("arb4j.version", "dev");
      var    lines = new java.util.ArrayList<String>();
      lines.add("pade (arb4j " + arb4j + ")");
      props.stringPropertyNames()
           .stream()
           .filter(k -> !k.equals("arb4j.version"))
           .sorted()
           .forEach(k -> lines.add("  " +
               k.replace(".version", "") + " " + props.getProperty(k)));
      return lines.toArray(String[]::new);
    }
  }
\end{lstlisting}

\textbf{Lines 111--135:} Reads \texttt{build-info.properties} (auto-generated by
\texttt{make}) and formats its contents as version strings. The first line is
\texttt{pade (arb4j <version>)} followed by every other property (gmp, mpfr,
flint, xdotool, swig, java versions) sorted alphabetically, with the
\texttt{.version} suffix stripped from the key names. Loads via
\texttt{try-with-resources} so the \texttt{InputStream} is always closed.

\section{Options Fields (Lines 137--207)}

\begin{lstlisting}[firstnumber=137]
@Option(names = "--mu", description = "fractional order \mu (default: 1)",
        defaultValue = "1")
String  mu;
...
@Option(names = "--N", description = "number of orthogonal polynomials (default: 10)",
        defaultValue = "10")
int     N;
...
@Option(names = "--bits", description = "working precision in bits (default: 128)",
        defaultValue = "128")
int     bits;
\end{lstlisting}

Every field annotated with \texttt{@Option} becomes a CLI flag. The complete
inventory:

\begin{tabularx}{\textwidth}{l l X}
\toprule
\textbf{Flag} & \textbf{Type} & \textbf{Purpose} \\
\midrule
\texttt{--mu} & String & Fractional order $\mu$ (parsed to \texttt{Real}) \\
\texttt{--P} & String & Coefficient $P(v)$ (default \texttt{"1"}) \\
\texttt{--Q} & String & Coefficient $Q(v)$ (default \texttt{"0"}) \\
\texttt{--R} & String & Coefficient $R(v)$ (default \texttt{"-1"}) \\
\texttt{--N} & int & Number of orthogonal polynomials (default 10) \\
\texttt{--bits} & int & Working precision (default 128) \\
\texttt{--no-reset-radius} & boolean & Don't zero ball radii \\
\texttt{--no-autoterm} & boolean & Don't auto-terminate at $0\in\Delta\beta_n$ \\
\texttt{--print-polynomials} & boolean & Print orthogonal polynomial tables \\
\texttt{--print-moments} & int & Print moment sequence ($-1$=off) \\
\texttt{--rough-heston} & boolean & Use rough Heston coefficients \\
\texttt{--lambda} & String & Rough Heston $\lambda$ (default \texttt{"0.3"}) \\
\texttt{--theta} & String & Rough Heston $\theta$ (default \texttt{"0.04"}) \\
\texttt{--nu} & String & Rough Heston $\nu$ (default \texttt{"0.3"}) \\
\texttt{--V0} & String & Rough Heston $V_0$ (default \texttt{"0.04"}) \\
\texttt{--rho} & String & Rough Heston $\rho$ (default \texttt{"-0.7"}) \\
\texttt{--T} & String & Rough Heston $T$ (default \texttt{"1.0"}) \\
\texttt{--expr} & String & Arbitrary expression to Pad\'e-resum \\
\texttt{--center} & String & Expansion center for \texttt{--expr} (default \texttt{"0"}) \\
\texttt{--plot} & String & Reference expression to plot \\
\texttt{--eval} & String & Comma-separated evaluation points \\
\texttt{--plot-left} & double & Left plot endpoint (default 0) \\
\texttt{--plot-right} & double & Right plot endpoint (default 1) \\
\texttt{--plot-points} & int & Plot sample points (default 200) \\
\bottomrule
\end{tabularx}

\section{Entry Point --- \texttt{main()} (Lines 209--213)}

\begin{lstlisting}[firstnumber=209]
public static void main(String[] args)
{
  int exitCode = new CommandLine(new PadePrinter()).execute(args);
  System.exit(exitCode);
}
\end{lstlisting}

\textbf{Line 211:} \texttt{picocli} instantiates \texttt{PadePrinter}, injects
the option values from \texttt{args}, then calls \texttt{run()} on it.
\texttt{System.exit} propagates the exit code (0 for success, non-zero for
errors).

\section{\texttt{run()} --- Mode Dispatch (Lines 216--222)}

\begin{lstlisting}[firstnumber=216]
public void run()
{
  if (expr != null)
  {
    runExpressionMode();
    return;
  }

  System.out.println("M\"untz--Pad\'e orthogonal polynomials ...");
  ...
\end{lstlisting}

\textbf{Lines 218--222:} If \texttt{--expr} was provided, dispatch to
\texttt{runExpressionMode()} which handles the arbitrary-expression path.
Otherwise fall through to the standard Riccati-equation path.

\newpage

\section{Standard Mode --- Header and Setup (Lines 224--237)}

\begin{lstlisting}[firstnumber=224]
System.out.println("M\"untz--Pad\'e orthogonal polynomials of the spectral-tau solution");
System.out.println("of the constant-coefficient (time-independent) fractional Riccati equation");
System.out.printf("  D^%s y(t) = P(v) + Q(v)\cdot y(t) + R(v)\cdot y(t)^2,  y(0) = 0%n", mu);
System.out.printf("Working precision: %d bits (~%d decimal digits)%n",
                   bits, (int) (bits * Math.log10(2)));
System.out.printf("First %d orthogonal polynomials in z = t^\mu%n", N);
System.out.println();

Real    \mu     = new Real().set(mu, bits);
Complex zeroV = new Complex().zero();

try ( RiccatiMuntzPadeFunctional eq = makeEquation(\mu))
{
  MuntzPadeApproximant approx =
      (MuntzPadeApproximant) eq.evaluate(zeroV, 1, bits, null);
  Context              ctx    = approx.context;
\end{lstlisting}

\textbf{Lines 224--229:} Print the banner. \texttt{bits * Math.log10(2)}
converts bits to approximate decimal digits.

\textbf{Line 231:} \texttt{mu} is parsed from its string form to a
\texttt{Real} ball at \texttt{bits} precision. The raw string \texttt{mu} came
from the \texttt{--mu} CLI option (a String, not a Real), because
\texttt{picocli} cannot directly inject into \texttt{arb.Real}.

\textbf{Line 232:} \texttt{zeroV = new Complex().zero()} creates a zero-valued
complex at the perturbation point $v = 0$. This is the evaluation point for the
$\sigma$-table: the orthogonal-polynomial coefficients are polynomials in $v$,
and evaluating at $v=0$ gives the scalar Pad\'e coefficients used in the
assembled rational approximant.

\textbf{Line 234:} \texttt{makeEquation($\mu$)} constructs the
\texttt{RiccatiMuntzPadeFunctional} (or
\texttt{RoughHestonRiccatiMuntzPadeFunctional} if \texttt{--rough-heston} is
set). The \texttt{try}-with-resources ensures the equation is closed when done.

\textbf{Line 236:} \texttt{eq.evaluate(zeroV, 1, bits, null)} evaluates the
functional at $v = 0$ with order 1. This triggers the entire M\"untz--Pad\'e
machinery: the $\sigma$-table sequences are compiled, the moment sequence is
computed from the Riccati recurrence, the Jacobi coefficients are extracted,
and the orthogonal polynomials are built. The return type is cast to
\texttt{MuntzPadeApproximant}, which holds all results.

\textbf{Line 237:} Extract the \texttt{Context} for later use in registering
function mappings for the plot/eval path.

\newpage

\section{Standard Mode --- Coefficient Functions (Lines 239--275)}

\begin{lstlisting}[firstnumber=239]
System.out.println("=".repeat(70));
System.out.println("Coefficient functions");
System.out.println("=".repeat(70));
System.out.println();
if (roughHeston)
{
  System.out.println("  P(v) = \frac12(-v^2 - iv)");
  ...
}
else
{
  System.out.printf("  P(v) = %s  (constant)%n", pCoeff);
  System.out.printf("  Q(v) = %s  (constant)%n", qCoeff);
  System.out.printf("  R(v) = %s  (constant)%n", rCoeff);
}
System.out.println();
eq.p.setIndependentVariableName("v");
eq.q.setIndependentVariableName("v");
eq.r.setIndependentVariableName("v");
System.out.printf("  p(v) = %s%n", eq.p);
System.out.printf("  q(v) = %s%n", eq.q);
System.out.printf("  r(v) = %s%n", eq.r);
System.out.println("=".repeat(70));
\end{lstlisting}

\textbf{Lines 239--267:} Print the coefficient functions $P$, $Q$, $R$. In
rough Heston mode, the coefficient formulas are hard-coded
($P(v) = \frac12(-v^2 - iv)$ etc.) because they are fixed by the rough Heston
model. Otherwise the user-provided strings from \texttt{--P}, \texttt{--Q},
\texttt{--R} are displayed.

\textbf{Lines 268--273:} Rename the independent variable of each compiled
coefficient expression to ``v'' for display, then print them as evaluated
polynomials. \texttt{eq.p}, \texttt{eq.q}, \texttt{eq.r} are
\texttt{ComplexPolynomial} instances or compiled expressions depending on the
implementation.

\section{Standard Mode --- $\sigma$-Table Definitions (Lines 277--301)}

\begin{lstlisting}[firstnumber=277]
System.out.println("=".repeat(70));
System.out.println("\sigma-table definitions (compiled expressions)");
System.out.println("=".repeat(70));
System.out.println();
try ( ComplexPolynomial disc = eq.discriminant(bits, new ComplexPolynomial()))
{
  disc.setIndependentVariableName("v");
  System.out.printf("  discriminant = Q^2 - 4PR = %s%n", disc);
}
System.out.println();
System.out.printf("  a(k) = %s%n", approx.aSeq.toNameString());
System.out.printf("  m(k) = %s%n", approx.mSeq.toNameString());
System.out.printf("  \sigma(j)(k) = %s%n", approx.\sigmaSeq.toNameString());
System.out.printf("  h(j) = %s%n", approx.hSeq.toNameString());
System.out.printf("  \alpha(j) = %s%n", approx.\alphaSeq.toNameString());
System.out.printf("  \beta(j) = %s%n", approx.\betaSeq.toNameString());
System.out.printf("  Pn(n) = %s%n", approx.PnSeq.toNameString());
System.out.println();
System.out.println("Pad\'e assembly (compiled expressions):");
System.out.printf("  \Phi_den(M) = %s%n", approx.\Phi denSeq.toNameString());
System.out.printf("  \Phi_num(M) = %s%n", approx.\Phi numSeq.toNameString());
System.out.printf("  \Phi(M)(z) = %s%n", approx.\Phi.toNameString());
System.out.printf("  K_n(z,w) = %s%n", approx.KnSeq.toNameString());
System.out.println("=".repeat(70));
\end{lstlisting}

\textbf{Line 281:} Compute and display the discriminant $Q^2 - 4PR$ as a
polynomial in $v$.

\textbf{Lines 287--293:} For each sequence in the approximant, call
\texttt{toNameString()} which returns the expression-compiler's source-level
representation of the compiled expression. This shows the user exactly how the
compiler transcribed the mathematical definitions into executable bytecode:

\begin{itemize}[nosep]
\item \texttt{a(k)} --- The Riccati coefficient sequence (moments before shift)
\item \texttt{m(k) = a(k+1)} --- The moment functional (1-index shift)
\item $\sigma(j)(k)$ --- The $\sigma$-table (Chebyshev--Stieltjes procedure)
\item $h(j) = \sigma(j)(j)$ --- Squared $L^2$ norms
\item $\alpha(j), \beta(j)$ --- Jacobi recurrence coefficients
\item $P_n(n)$ --- Associated (second-kind) polynomials
\item $\Phi_{\text{den}}(M)$ --- Denominator polynomial (monic orthogonal, favard recurrence)
\item $\Phi_{\text{num}}(M)$ --- Numerator polynomial (associated)
\item $\Phi(M)(z)$ --- The assembled Pad\'e approximant $\Phi_{\text{num}} / \Phi_{\text{den}}$
\item $K_n(z,w)$ --- The Christoffel--Darboux reproducing kernel
\end{itemize}

\section{Standard Mode --- Moment Sequence Printing (Lines 303--325)}

\begin{lstlisting}[firstnumber=303]
if (printMoments >= 0)
{
  int numMoments = printMoments == 0 ? N : printMoments;
  System.out.println("=".repeat(70));
  System.out.printf("Moment sequence m(k) = a(k+1)  [first %d moments]%n", numMoments);
  ...
  for (int k = 0; k < numMoments; k++)
  {
    moments[k][0] = String.valueOf(k);
    try ( var mk = approx.ops.m.evaluate(k, bits))
    {
      mk.printPrecision = true;
      moments[k][1] = mk.toString();
    }
  }
  new TextTable(cols, moments).printTable();
}
\end{lstlisting}

\textbf{Lines 303--304:} Only print moments when \texttt{--print-moments >= 0}.
\texttt{printMoments == 0} means ``print $N$ moments'' (same as the polynomial
count). A positive value gives an explicit count.

\textbf{Line 316:} \texttt{approx.ops.m.evaluate(k, bits)} evaluates the $k$-th
moment $m(k) = a(k+1)$. The moment is a \texttt{ComplexPolynomial} in $v$,
evaluated at the frozen point $v = 0$ inside the approximant construction, but
here we display it as a polynomial.

\textbf{Line 318:} \texttt{mk.printPrecision = true} enables printing the ball
radius alongside the midpoint, revealing the certified error bound.

\section{Standard Mode --- Polynomial Tables (Lines 327--387)}

\begin{lstlisting}[firstnumber=327]
ComplexPolynomialSequence PnSeq = approx.PnSeq;
ComplexSequence           \alpha vSeq = approx.\alpha vSeq;
ComplexSequence           \beta vSeq = approx.\beta vSeq;
ComplexSequence           hvSeq = approx.hvSeq;
\end{lstlisting}

\textbf{Lines 327--330:} Local aliases for the scalar-evaluated sequences.
These are the $\alpha$, $\beta$, $h$ values evaluated at $v=0$ (the
perturbation point), giving scalar (complex-number) coefficients rather than
polynomials in $v$.

\newpage

\begin{lstlisting}[firstnumber=332]
if (printPolynomials)
{
  // Denominator polynomials P_n(z) -- monic orthogonal, first kind
  ...
  for (int n = 0; n < N; n++)
  {
    try ( ComplexPolynomial Pn = approx.ops.evaluate(n, bits))
    {
      Pn.setIndependentVariableName("z");
      ...
    }
  }
  new TextTable(polyCols, polyData).printTable();

  // Numerator polynomials Pn_n(z) -- associated, second kind
  ...
  for (int n = 0; n < N; n++)
  {
    try ( ComplexPolynomial PnAssoc = PnSeq.evaluate(n, bits))
    {
      PnAssoc.setIndependentVariableName("z");
      ...
    }
  }
  new TextTable(polyCols, assocData).printTable();
}
\end{lstlisting}

\textbf{Lines 332--387:} Conditionally print both sets of orthogonal
polynomials:

\begin{enumerate}[nosep]
\item \textbf{Denominator polynomials $P_n(z)$} (first kind, monic orthogonal):
  Evaluated via \texttt{approx.ops.evaluate(n, bits)} which returns a
  \texttt{ComplexPolynomial}. These are the orthogonal polynomials with respect
  to the moment functional, satisfying the three-term recurrence
  $P_{n+1}(z) = (z - \alpha_n) P_n(z) - \beta_n P_{n-1}(z)$.

\item \textbf{Numerator polynomials $\tilde P_n(z)$} (second kind, associated):
  Evaluated via \texttt{PnSeq.evaluate(n, bits)}. These satisfy the same
  recurrence but with different initial conditions.
\end{enumerate}

Each polynomial's independent variable is renamed to ``z'' for display.

\section{Standard Mode --- Jacobi Coefficient Table (Lines 389--449)}

\begin{lstlisting}[firstnumber=394]
String[]   jacobiCols = { "n", "\alpha_n", "\beta_n", "h_n", "\Delta\beta_n",
                          "0\in\Delta\beta_n" };
int        actualN    = N;
String[][] jacobiData = new String[N][6];
try ( Complex \alpha n = new Complex(); Complex \beta n = new Complex();
      Complex hn = new Complex(); Complex \beta nPrev = new Complex();
      Complex delta\beta = new Complex())
{
  for (int n = 0; n < N; n++)
  {
    \alpha vSeq.evaluate(n, bits, \alpha n);
    \beta vSeq.evaluate(n, bits, \beta n);
    hvSeq.evaluate(n, bits, hn);
    ...
\end{lstlisting}

This is the computational core of the display. For each $n = 0, \ldots, N-1$:

\textbf{Lines 402--404:} Evaluate the three Jacobi sequences at index $n$,
writing into pre-allocated \texttt{Complex} receivers (the ``last-argument-is-
result'' pattern). Reusing the same five \texttt{Complex} instances across all
iterations avoids allocation churn.

\textbf{Lines 405--413:} If \texttt{--no-reset-radius} is NOT set (the default),
zero out the ball radii of the results. The $\sigma$-table's ball arithmetic
propagates uncertainty from the moment sequence through the recurrence. For
display purposes, we want to see the value without the radius clouding the
output. The radii are still correct---they just aren't printed.

\textbf{Lines 414--420:} Check for NaN. Ball arithmetic produces NaN when
the computation loses all precision (e.g., catastrophic cancellation). If
detected, terminate the table early.

\textbf{Lines 425--443:} For $n \ge 1$, compute $\Delta\beta_n = \beta_n - \beta_{n-1}$.
\texttt{deltaBeta.getReal().containsZero() \&\& deltaBeta.getImag().containsZero()}
checks whether the ball interval for $\Delta\beta_n$ contains zero. If it does
and \texttt{--no-autoterm} is not set, the iteration auto-terminates: the
$\sigma$-table has converged and further rows would not change the recurrence
coefficients.

\textbf{Line 446--448:} Trim the table to the actual number of rows that were
computed and print it via \texttt{TextTable}.

\section{Standard Mode --- Plotting and Evaluation (Lines 451--507)}

\begin{lstlisting}[firstnumber=451]
if (plotExpr != null || evalPoints != null)
{
  int M = actualN - 1;
  try ( Integer Mi = new Integer(M);
        ComplexPolynomial numPoly = approx.\Phi numSeq.evaluate(Mi, bits);
        ComplexPolynomial denPoly = approx.\Phi denSeq.evaluate(Mi, bits))
  {
    numPoly.setIndependentVariableName("z");
    denPoly.setIndependentVariableName("z");
    ctx.registerFunctionMapping("num_fn", numPoly, Complex.class, Complex.class);
    ctx.registerFunctionMapping("den_fn", denPoly, Complex.class, Complex.class);
    RealFunction padeOfT = RealFunction.express(
        "t\gg re(num_fn(t^" + mu + ")/den_fn(t^" + mu + "))", ctx);
    ...
  }
}
\end{lstlisting}

\textbf{Line 453:} The Pad\'e order $M$ is one less than the number of
Jacobi rows (the last computed $n$).

\textbf{Lines 454--456:} Evaluate the numerator and denominator polynomials of
the Pad\'e approximant at order $M$. These are \texttt{ComplexPolynomial}
instances with coefficients that are balls---the certified error bounds.

\textbf{Lines 460--462:} Register the numerator and denominator as functions
\texttt{num\_fn} and \texttt{den\_fn} in the \texttt{Context}, then compile a
real-valued expression \texttt{t $\rightarrow$ re(num\_fn(t\textasciicircum mu)/den\_fn(t\textasciicircum mu))}
that evaluates the Pad\'e approximant at any point $t$. The change of variable
$z = t^{\mu}$ is composed into the expression.

\textbf{Line 464--468 (plot):} If \texttt{--plot} was given, compile the
reference expression and call \texttt{ShellFunctions.plot()} which opens a
JavaFX window showing both functions.

\textbf{Lines 471--501 (eval):} If \texttt{--eval} was given, iterate over the
comma-separated points, evaluate both the Pad\'e approximant and the reference
(if any) at each point, compute the absolute difference, and display as a
table. Each evaluation is a ball-arithmetic computation that carries a
certified radius.

\section{Expression Mode --- \texttt{runExpressionMode()} (Lines 511--750)}

\begin{lstlisting}[firstnumber=511]
private void runExpressionMode()
{
  System.out.println("Pad\'e approximant of arbitrary expression");
  System.out.printf("  f(t) = %s%n", expr);
  System.out.printf("  expansion center: t = %s%n", center);
  ...
  try ( ComplexFunction f = ComplexFunction.express("t\gg" + expr))
  {
    int            numMoments = N + 2;
    ComplexFunction compiled   = f;
    try ( Complex t0 = new Complex(new Real(center, bits));
          Complex taylorCoefficients = Complex.newVector(numMoments))
    {
      taylorCoefficients.zero();
      compiled.evaluate(t0, numMoments, bits, taylorCoefficients);
      ...
\end{lstlisting}

\textbf{Line 520:} Compile the user's expression \texttt{expr} into a
\texttt{ComplexFunction}. The string \texttt{"t $\rightarrow$ " + expr} parses
the expression compiler's Unicode syntax, builds an AST, and emits JVM bytecode.

\textbf{Line 522:} Request $N+2$ Taylor coefficients (slightly more than the
number of orthogonal polynomials to give the $\sigma$-table enough moments).

\textbf{Line 524:} Build the expansion center $t_0$ as a \texttt{Complex} ball.

\textbf{Line 525:} Allocate a vector of $N+2$ complex numbers to hold the
Taylor coefficients.

\textbf{Line 528:} Evaluate the compiled function at $t_0$ with order $N+2$.
The expression compiler's jet mode computes all derivatives simultaneously,
producing $c_k = f^{(k)}(t_0)/k!$ in one pass.

\newpage

\begin{lstlisting}[firstnumber=546]
try ( MuntzPadeApproximant approx =
          makeExpressionApproximant(taylorCoefficients))
{
  Context ctx = approx.context;
  ...
  // Rest is identical to standard mode:
  //   - Print \sigma-table definitions
  //   - Print polynomial tables (if --print-polynomials)
  //   - Print Jacobi coefficient table with auto-termination
  //   - Plot/eval (if --plot or --eval)
\end{lstlisting}

\textbf{Line 546:} Construct a \texttt{MuntzPadeApproximant} from the Taylor
coefficients via \texttt{makeExpressionApproximant}. From this point onward,
the expression mode reuses the same logic as the standard mode: printing the
$\sigma$-table definitions, polynomial tables, Jacobi coefficients, and
plot/eval sections. The code structure is identical (\texttt{approx.$\sigma$Seq},
\texttt{approx.hSeq}, etc.) because the approximant returns the same typed
sequences regardless of whether the moments came from the Riccati recurrence
or from Taylor coefficients of an arbitrary expression.

\section{\texttt{makeExpressionApproximant()} (Lines 752--774)}

\begin{lstlisting}[firstnumber=752]
private MuntzPadeApproximant makeExpressionApproximant(
    Complex taylorCoefficients)
{
  Real    \alpha = new Real("1", bits);
  Complex v      = new Complex().zero();

  int numCoefficients = taylorCoefficients.dim();

  ComplexPolynomialSequence a = (n, order, bits2, result) ->
  {
    int k = n.getSignedValue();
    if (k >= 1 && k <= numCoefficients)
    {
      result.set(taylorCoefficients.get(k - 1));
    }
    else
    {
      result.zero();
    }
    return result;
  };

  return new MuntzPadeApproximant(\alpha, a, v, bits);
}
\end{grammar}

\textbf{Line 754:} The fractional order $\alpha$ is set to 1 (ordinary Pad\'e,
not M\"untz--Pad\'e), because the expression mode works with standard Taylor
series $z = t$ rather than $z = t^{\mu}$.

\textbf{Line 755:} The perturbation point $v$ is set to 0 (no parameter
dependence).

\textbf{Lines 759--771:} Define the coefficient sequence $a(k)$ as a lambda that
returns the $k$-th Taylor coefficient for $1 \le k \le N$, and zero otherwise.
The 1-index shift matches the convention of the Riccati recurrence where
$a(1)$ is the first non-trivial coefficient.

\textbf{Line 773:} Construct and return the \texttt{MuntzPadeApproximant}. This
triggers the same $\sigma$-table compilation, Jacobi coefficient extraction,
and orthogonal polynomial construction as the Riccati path---all driven by the
moment sequence $m(k) = a(k+1)$.

\section{\texttt{makeEquation()} (Lines 776--809)}

\begin{lstlisting}[firstnumber=776]
private RiccatiMuntzPadeFunctional makeEquation(Real \mu)
{
  if (roughHeston)
  {
    Context ctx  = new Context();
    Real    \lambda    = new Real(\lambda, bits).setName("\lambda");
    Real    \theta    = new Real(\theta, bits).setName("\theta");
    Real    \nu       = new Real(nu, bits).setName("\nu");
    Real    V0r       = new Real(V0, bits).setName("V0");
    Real    \rho      = new Real(rho, bits).setName("\rho");
    Real    Tval      = new Real(T, bits).setName("T");
    ctx.registerVariable(\lambda);
    ctx.registerVariable(\theta);
    ctx.registerVariable(\nu);
    ctx.registerVariable(V0r);
    ctx.registerVariable(\rho);
    ctx.registerVariable(Tval);
    return new RoughHestonRiccatiMuntzPadeFunctional(ctx, \mu);
  }
  else
  {
    return new RiccatiMuntzPadeFunctional(\mu, pCoeff, qCoeff, rCoeff);
  }
}
\end{lstlisting}

\textbf{Lines 778--801 (rough Heston path):} Creates a \texttt{Context} and
registers each rough Heston parameter as a named variable. The
\texttt{RoughHestonRiccatiMuntzPadeFunctional} uses these variables in its
compiled coefficient expressions, which are strings like
\texttt{"(1/2)*(-v$^\wedge$2 - i*v)"} for $P(v)$. The expression compiler
resolves the parameter names ($\lambda$, $\theta$, $\nu$, etc.) at compile
time through the context.

\textbf{Lines 802--808 (standard path):} Constructs a
\texttt{RiccatiMuntzPadeFunctional} directly from the user-provided strings
\texttt{pCoeff}, \texttt{qCoeff}, \texttt{rCoeff}, and the fractional order
$\mu$. No context is needed because the coefficients are literal polynomial
strings.

\section{Data Flow Summary}

\begin{enumerate}
\item \textbf{Input parsing:} \texttt{picocli} injects CLI flags into fields.
\item \textbf{Equation construction:} \texttt{makeEquation($\mu$)} builds the
  Riccati functional, compiling $P$, $Q$, $R$ into executable bytecode.
\item \textbf{Evaluation:} \texttt{eq.evaluate(zeroV, 1, bits, null)} triggers
  the full M\"untz--Pad\'e pipeline:
  \begin{enumerate}
  \item The Riccati recurrence generates the coefficient sequence $a(k)$.
  \item $m(k) = a(k+1)$ defines the moment functional.
  \item The $\sigma$-table (Chebyshev--Stieltjes procedure) reduces the Hankel
    moment matrix to tridiagonal form, producing $\alpha(j)$, $\beta(j)$, $h(j)$.
  \item The orthogonal polynomials $P_n(z)$ (first kind) and associated
    polynomials $\tilde P_n(z)$ (second kind) are built via Favard's recurrence.
  \item The Pad\'e numerator $\Phi_{\text{num}}(M)$ and denominator
    $\Phi_{\text{den}}(M)$ are assembled.
  \item The Christoffel--Darboux kernel $K_n(z,w)$ is compiled.
  \end{enumerate}
\item \textbf{Display:} The printer formats and prints every intermediate and
  final result as human-readable tables.
\item \textbf{Termination:} The $\sigma$-table auto-terminates when
  $0 \in \Delta\beta_n$ (the ball arithmetic detects convergence).
\item \textbf{Plot/Evaluate:} The assembled rational approximant is registered
  as a context function and used for point evaluation or plotting against a
  reference expression.
\end{enumerate}

\section{Design Notes}

\subsection{Why \texttt{System.out.println} Instead of SLF4J?}

PadePrinter is a CLI application (not a library or test). Its entire purpose
is to write formatted output to stdout for the user to read. Using
\texttt{System.out} here is the correct choice---SLF4J is for programmatic
logging, not user-facing output.

\subsection{Ball Arithmetic and the Radius Reset}

The \texttt{--no-reset-radius} flag controls whether the ball radii are zeroed
for display. The radii are mathematically correct and propagate all
uncertainties from the moment sequence through the $\sigma$-table. Zeroing them
gives a cleaner display of the midpoint values without hiding the fact that the
radii were present. The default behavior resets them; passing
\texttt{--no-reset-radius} preserves them so the user can inspect the error
enclosures.

\subsection{Auto-Termination by $0 \in \Delta\beta_n$}

The $\sigma$-table's convergence criterion is exact, not heuristic. When the
moment sequence is exhausted (the Jacobi recurrence reaches a fixed point),
$\beta_n = \beta_{n-1}$ and therefore $\Delta\beta_n = 0$. Ball arithmetic
detects this via \texttt{containsZero()} on both the real and imaginary
components. This is the same termination condition used in the computed
$\sigma$-table theory (see \texttt{docs/RKHS-Convergence-Framework.md}).

\subsection{Resource Management}

Every native-memory-backed type (\texttt{Real}, \texttt{Complex},
\texttt{ComplexPolynomial}, \texttt{MuntzPadeApproximant}) is either wrapped in
\texttt{try}-with-resources or has its ownership transferred to a
\texttt{try}-with-resources scope. The \texttt{RiccatiMuntzPadeFunctional} and
its contained \texttt{MuntzPadeApproximant} form a tree of closeable objects;
the outer \texttt{try} block ensures the root is closed, which cascades through
the tree.

\subsection{Two-Mode Code Duplication}

The standard mode (Riccati equation) and expression mode (arbitrary function)
share approximately 250 lines of identical logic for printing the
$\sigma$-table, polynomial tables, Jacobi coefficients, and plot/eval sections.
This duplication is structural rather than accidental: the two modes construct
the \texttt{MuntzPadeApproximant} through different paths (Riccati recurrence
vs. Taylor coefficient extraction), but once the approximant exists, the
display logic is identical. A future refactoring could extract the common
display code into shared methods.

\end{document}
